% ============================================================================
%  Part VIII: Comparison with Existing Formats
% ============================================================================
\section{Comparison with Existing Formats}
\label{sec:compare}

PQID is one option among many. This part walks through the formats most likely to be on the table --- UUIDv4, UUIDv7, ULID, KSUID, and Snowflake --- and explains where PQID is the better choice and where the existing format is the right answer. The intent is to make the tradeoffs explicit, not to claim that PQID is universally superior.

\subsection{vs UUIDv4 \citep{rfc4122}}

UUIDv4 is the universal default, supported in every language's standard library and every database engine. It is 16 bytes, 22 characters Base64URL or 36 characters hyphenated. Its 122 random bits give 61-bit post-quantum security.

\textbf{PQID wins when:} time-sortability matters (UUIDv4 is randomly ordered, which destroys B-tree performance), or post-quantum margins matter (61 bits is on the edge), or you want explicit distributed-systems support (UUIDv4 has none).

\textbf{UUIDv4 wins when:} byte cost dominates (UUIDv4 is $2.5\times$ smaller), the existing tooling ecosystem matters (every language's standard library handles UUIDs natively), or you genuinely don't care about ordering (analytical workloads with shuffled access patterns are an example).

\subsection{vs UUIDv7 \citep{rfc9562} --- the closest competitor, in detail}

UUIDv7 is the IETF's 2024 response to the demand for time-ordered UUIDs. Of all the formats discussed in this paper, it is the one most likely to be on the table when PQID is being evaluated, because it is the only RFC-standardized time-sortable UUID variant and because it is the format that ships in the language standard libraries published in 2024 and 2025. It deserves a careful comparison rather than a one-line dismissal.

\subsubsection{Layout and entropy}

UUIDv7 packs into 16 bytes: 48 bits of millisecond timestamp, 4 bits fixed to the value 7 (version), 12 bits of ``rand\_a,'' 2 bits fixed to UUID's standard variant marker, and 62 bits of ``rand\_b'' --- 74 random bits total, distributed across two non-contiguous fields. Compared to PQID:

\begin{center}
\renewcommand{\arraystretch}{1.15}
\begin{tabular}{lll}
\toprule
\textbf{Property} & \textbf{UUIDv7} & \textbf{PQID} \\
\midrule
Total size                  & 16 bytes / 22 B64 chars       & 41 bytes / 55 B64 chars \\
Random bits                 & 74                            & 256 \\
Classical PQ security       & 37 bits                       & 128 bits \\
Quantum PQ security (Grover)& $\sim$37 bits                 & $\sim$128 bits \\
Intra-ms monotonic order    & Implementation-defined (RFC §6.2) & Guaranteed (counter) \\
Distributed-systems variant & None (wall-clock only)        & HLC (variant 0x02) \\
\bottomrule
\end{tabular}
\end{center}

\subsubsection{The intra-millisecond ordering footgun}

RFC 9562 §6.2 lays out three permissible methods for generating UUIDv7s with monotonic intra-ms ordering. Method 1 (use the 12 \code{rand\_a} bits as a counter) and Method 3 (use 12 bits of monotonic counter that re-randomizes per ms) provide intra-ms monotonicity at a cost: Method 1 gives only $4{,}096$ counter values per ms, and Method 3's re-randomization makes verification harder. Method 2 (just use random bits with no counter) is permitted by the spec and produces UUIDs that are NOT monotonically ordered within a single millisecond. Two consecutive calls falling in ms $T$ produce IDs ordered randomly relative to each other. The RFC explicitly permits this.

In benchmarks against the Python \code{uuid} stdlib (which uses a Method-1-like approach), UUIDv7 sortedness measured 50.18\% --- statistically indistinguishable from random ordering of consecutive IDs. PostgreSQL's \code{uuidv7()} function (added in v18) uses Method 3 and produces fully-sorted IDs. Java's \code{java-uuid-generator} library defaults closer to Method 1. Go's standard library does not have v7 at all (third-party only). The implementation choice is invisible from the wire format: a UUIDv7 produced by Method 1 looks identical to one produced by Method 3, and a system that mixes generators across languages cannot recover the ordering guarantee from inspection alone.

\begin{calloutwarning}[Why this matters]
Time-sortable identifier formats deliver their value (B-tree locality, replication ordering, log analysis) through the assumption that sorting reflects generation time. UUIDv7 partially honors this assumption between milliseconds but only conditionally honors it within milliseconds, and the conditionality is not visible from the format. In a heterogeneous deployment that spans multiple languages, frameworks, or library versions --- exactly the deployment shape of large clusters --- the assumption holds for some IDs and fails for others, with no way to distinguish them downstream. PQID makes the guarantee unconditional.
\end{calloutwarning}

\subsubsection{Distributed-systems support}

UUIDv7's timestamp is wall-clock milliseconds since the Unix epoch. There is no provision for logical-clock semantics, no field for capturing causal observations from peers, and no drift bound. Two UUIDv7s generated on different nodes whose clocks differ by 50 ms --- a routine cross-data-center NTP skew --- sort according to wall-clock time, which may or may not reflect causal time.

This is fine for the use cases UUIDv7 was actually designed for: request IDs in a stateless service, primary keys in a single database, span IDs in a distributed trace where the trace itself records causal relationships separately. It is not fine for use cases that need the identifier itself to encode causal ordering: distributed event logs spanning multiple nodes, CRDT version vectors, gradient-update streams in distributed training, multi-agent interaction traces. These cases need HLC semantics, which UUIDv7 cannot provide and PQID's variant 0x02 makes a first-class option.

\subsubsection{Post-quantum margins for long-lived identifiers}

UUIDv7's 74-bit random field gives 37 bits of effective security against Grover-style quantum search for a targeted second-preimage. For ephemeral identifiers (request IDs measured in seconds, span IDs measured in milliseconds), 37 bits is more than enough --- no plausible attacker has a quantum computer that can find a specific 74-bit value in any window where the identifier remains operationally meaningful. For long-lived identifiers (training-corpus records, audit logs, archival keys, financial transaction identifiers in long-retention systems), 37 bits is genuinely thin in a world where quantum-relevant attackers may exist by 2045--2050. PQID's 128-bit post-quantum margin is calibrated to the same long-horizon threat model that drives NIST's selection of 128-bit symmetric security as the appropriate level through 2030 and beyond.

Pragmatically, this means UUIDv7 is a good choice for transient and operational identifiers, and a worse choice for identifiers that may need to remain unforgeable for decades. The two regimes coexist in most systems --- a frontier-AI training run produces both ephemeral request IDs (UUIDv7 fine) and long-archived training-corpus records (UUIDv7 risky). PQID's properties make it suitable for both regimes; the cost is the byte overhead.

\subsubsection{When UUIDv7 is the right answer}

Standards compliance matters (UUIDv7 is an RFC, PQID is not yet); the 16-byte fixed size matters; the existing UUID tooling ecosystem (database UUID column types, language UUID libraries, tracing systems built around UUID conventions) provides operational value worth preserving; the use case is operational/transient rather than archival; the deployment is single-tenant with controlled implementation choice; or causal ordering across nodes is genuinely not part of the requirements. For these cases, UUIDv7 is well-designed and PQID is overkill.

When PQID is the right answer: post-quantum margins matter for long-horizon identifier defensibility; intra-millisecond ordering must be guaranteed across heterogeneous implementations; causal ordering across distributed nodes is part of the requirements (HLC variant); or the deployment scale is large enough that the small per-format failures in UUIDv7 (Method-1 sortedness, 37-bit PQ margin) compound into observable defects. Part \ref{sec:ai-scale} works through the specific case of million-node clusters in detail.

\subsection{vs ULID \citep{ulid}}

ULID is a popular community standard: 48 bits of ms timestamp + 80 bits of randomness, 16 bytes total. ULID guarantees lexicographic time ordering and provides a Crockford Base32 encoding.

\textbf{PQID wins when:} post-quantum margins matter (ULID's 40 post-quantum bits is too thin for new designs targeting 2030+), or distributed-systems semantics matter (ULID is wall-clock only). PQID's design is essentially ``ULID with more entropy and a distributed-systems variant''; if you like ULID, PQID is the natural upgrade path.

\textbf{ULID wins when:} byte cost matters (ULID is 39\% smaller than PQID), the broad ecosystem support matters (ULID has libraries in roughly every language now), or post-quantum security is genuinely not a concern for your use case.

\subsection{vs KSUID \citep{ksuid}}

KSUID: 32-bit second timestamp + 128 random bits, 20 bytes total. Time-ordered, but only at second granularity, which means many KSUIDs generated within a single second are ordered randomly relative to each other.

\textbf{PQID wins when:} sub-second ordering matters (KSUID does not provide it), post-quantum margins matter (KSUID's 64 bits is on the edge), or distributed-systems semantics matter.

\textbf{KSUID wins when:} byte cost matters and you want the largest random field possible at that size (KSUID's 128-bit random field is the largest in the under-32-byte category), or you happen to be on a stack where Segment's KSUID library is already deeply embedded.

\subsection{vs Snowflake \citep{snowflake}}

Twitter's Snowflake: 41-bit ms timestamp + 10-bit machine ID + 12-bit sequence, packed into a 64-bit integer. Eight bytes total, with the smallest possible representation. Requires coordinated machine-ID assignment.

\textbf{PQID wins when:} you can't coordinate machine-ID assignment (PQID is uncoordinated; Snowflake requires every generator to have a unique pre-assigned ID), or you need cryptographic-grade uniqueness (Snowflake has zero random bits and zero collision resistance --- its uniqueness depends entirely on machine-ID assignment correctness), or distributed-systems semantics matter.

\textbf{Snowflake wins when:} bytes are at an absolute premium (8 vs 41 is $5\times$ smaller), and you have a coordination service (e.g., Zookeeper) that can hand out unique machine IDs. For purely intra-database identifiers in a coordinated cluster, Snowflake's compactness is unbeatable.

\subsection{Summary table}

\begin{center}
\renewcommand{\arraystretch}{1.15}
\footnotesize
\begin{tabular}{lcccccll}
\toprule
\textbf{Concern} & \textbf{UUIDv4} & \textbf{UUIDv7} & \textbf{ULID} & \textbf{KSUID} & \textbf{Snowflake} & \textbf{PQID Std} & \textbf{PQID HLC} \\
\midrule
Sortable           & ---       & Partial    & Yes (ms)   & Yes (s)    & Yes        & \textbf{Yes (ms)}   & \textbf{Yes (HLC)} \\
PQ ($\geq$ 128 b)  & ---       & ---        & ---        & ---        & ---        & \textbf{Yes}        & \textbf{Yes} \\
Distributed/HLC    & ---       & ---        & ---        & ---        & Coord.     & ---                 & \textbf{Yes} \\
Bytes              & 16        & 16         & 16         & 20         & 8          & 41                  & 41 \\
B64 length         & 22        & 22         & 22         & 27         & 11         & 55                  & 55 \\
RFC / standard     & RFC       & RFC        & Community  & Community  & Industry   & Proposed            & Proposed \\
\bottomrule
\end{tabular}
\end{center}

